% ---- score2 standard preamble (from: python3 score2_tools.py preamble) ----
\documentclass[10pt]{extarticle}
\usepackage[a4paper,margin=1.5cm]{geometry}
\usepackage[T1]{fontenc}
\usepackage{amsmath,amssymb}
\usepackage{graphicx}
\usepackage[export]{adjustbox}
\usepackage{multicol}
\usepackage{enumitem}
\usepackage{booktabs,array,tabularx}
\setlength{\columnsep}{0.7cm}
\setlength{\parindent}{0pt}
\setlength{\parskip}{3pt}
\pagestyle{empty}
\makeatletter
% figure registry (figdefs.tex is generated by the extract step)
\newcommand{\qfigdef}[4]{\@namedef{qf@file@#1}{#2}\@namedef{qf@w@#1}{#3}\@namedef{qf@wide@#1}{#4}}
\newcommand{\qf@check}[1]{\@ifundefined{qf@file@#1}{\PackageError{score2}{Unknown figure id #1}{}}{}}
\newcommand{\qf@img}[4][]{\includegraphics[width=#2\dimexpr\@nameuse{qf@w@#4}pt\relax,max width=#3,max height=0.45\textheight,#1]{\@nameuse{qf@file@#4}}}
% Full-width questions: finalize wraps every question that contains a wide figure in \qfullbegin...\qfullend,
% so the WHOLE question (stem, figures, options) is set across the page in reading order, then the two columns
% resume. The question's head - its number, stem and first figure (\qfullsplit marks the end of the head) -
% is kept on one page, so a number is never left at a page foot with its figure overleaf; the rest may break
% across pages normally, so no big gaps are left. Without a \qfullsplit, the whole question is kept together
% if it is under 45% of a page. Never write these yourself.
\newif\ifqf@full
\newif\ifqf@boxed
\newsavebox{\qf@box}
\newcommand{\qf@open}{\global\qf@fulltrue\par\global\qf@boxedtrue
  \setbox\qf@box\vbox\bgroup\hsize\textwidth\linewidth\textwidth}
\newcommand{\qf@emit}[1]{\par\egroup\global\qf@boxedfalse
  \ifdim\dimexpr\ht\qf@box+\dp\qf@box\relax<#1\textheight\noindent\usebox{\qf@box}\par\else\unvbox\qf@box\fi}
\newcommand{\qfullsplit}{\ifqf@boxed\qf@emit{0.85}\fi}
\newcommand{\qf@close}{\ifqf@boxed\qf@emit{0.45}\fi\global\qf@fullfalse}
\newcommand{\qfullbegin}{\end{multicols}\qf@open}
\newcommand{\qfullend}{\qf@close\begin{multicols}{2}}
\newcommand{\qfullstart}{\qf@open}
\newcommand{\qfulldone}{\qf@close}
% \qfig[scale]{id}: block figure, centred, directly where it is written (size set by the tool).
% Fallback if a wide figure is not inside a full-width question: break out of the columns in place - no floats.
\newcommand{\qfig}[2][1]{\qf@check{#2}%
  \ifqf@full\qf@blk{#1}{#2}\else
  \ifnum\@nameuse{qf@wide@#2}=1 \expandafter\@firstoftwo\else\expandafter\@secondoftwo\fi
  {\end{multicols}\begin{center}\qf@img{#1}{\textwidth}{#2}\end{center}\begin{multicols}{2}}%
  {\qf@blk{#1}{#2}}\fi}
\newcommand{\qf@blk}[2]{\par\smallskip{\centering\qf@img{#1}{\linewidth}{#2}\par}\smallskip}
% \qfigi[scale]{id}: inline figure (inside option lists / tables / minipages); never breaks columns.
\newcommand{\qfigi}[2][1]{\qf@check{#2}\qf@img[valign=c]{#1}{\linewidth}{#2}}
\makeatother
% structure macros - lint reads these to check completeness and figure placement
\newcommand{\papersec}[1]{\par\bigskip{\centering\Large\bfseries PAPER-#1\par}\medskip}
\newcommand{\partsec}[2]{\par\medskip\noindent{\large\bfseries PART-#1: #2}\par\smallskip}
\newcommand{\sect}[1]{\par\smallskip\noindent\textbf{#1}\par}
\newcommand{\qn}[1]{\par\smallskip\noindent\textbf{#1.}~}
\newcommand{\paraq}[2]{\par\smallskip\noindent\textbf{Paragraph for Questions #1 and #2}\par}
\newenvironment{qoptions}{\begin{itemize}[leftmargin=1.8em,itemsep=1pt,topsep=2pt,label={}]}{\end{itemize}}
\newcommand{\opt}[2]{\item[(#1)] #2}
% figure registry (generated by score2_tools.py finalize)
\qfigdef{p1_phy_s1_q01_1}{Score-II - Test 34 (02 Feb 2025)-img/p1_phy_s1_q01_1.png}{106.9}{0}
\qfigdef{p1_phy_s1_q02_1}{Score-II - Test 34 (02 Feb 2025)-img/p1_phy_s1_q02_1.png}{86.3}{0}
\qfigdef{p1_phy_s1_q02_A}{Score-II - Test 34 (02 Feb 2025)-img/p1_phy_s1_q02_A.png}{81.2}{0}
\qfigdef{p1_phy_s1_q02_B}{Score-II - Test 34 (02 Feb 2025)-img/p1_phy_s1_q02_B.png}{75.9}{0}
\qfigdef{p1_phy_s1_q02_C}{Score-II - Test 34 (02 Feb 2025)-img/p1_phy_s1_q02_C.png}{81.0}{0}
\qfigdef{p1_phy_s1_q02_D}{Score-II - Test 34 (02 Feb 2025)-img/p1_phy_s1_q02_D.png}{81.2}{0}
\qfigdef{p1_phy_s1_q04_1}{Score-II - Test 34 (02 Feb 2025)-img/p1_phy_s1_q04_1.png}{96.5}{0}
\qfigdef{p1_phy_s2_q07_1}{Score-II - Test 34 (02 Feb 2025)-img/p1_phy_s2_q07_1.png}{130.7}{0}
\qfigdef{p1_phy_s2_q09_1}{Score-II - Test 34 (02 Feb 2025)-img/p1_phy_s2_q09_1.png}{87.6}{0}
\qfigdef{p1_phy_s2_q09_5}{Score-II - Test 34 (02 Feb 2025)-img/p1_phy_s2_q09_5.png}{81.2}{0}
\qfigdef{p1_phy_s2_q09_2}{Score-II - Test 34 (02 Feb 2025)-img/p1_phy_s2_q09_2.png}{86.8}{0}
\qfigdef{p1_phy_s2_q09_6}{Score-II - Test 34 (02 Feb 2025)-img/p1_phy_s2_q09_6.png}{77.1}{0}
\qfigdef{p1_phy_s2_q09_3}{Score-II - Test 34 (02 Feb 2025)-img/p1_phy_s2_q09_3.png}{92.1}{0}
\qfigdef{p1_phy_s2_q09_7}{Score-II - Test 34 (02 Feb 2025)-img/p1_phy_s2_q09_7.png}{76.9}{0}
\qfigdef{p1_phy_s2_q09_4}{Score-II - Test 34 (02 Feb 2025)-img/p1_phy_s2_q09_4.png}{92.1}{0}
\qfigdef{p1_phy_s2_q09_8}{Score-II - Test 34 (02 Feb 2025)-img/p1_phy_s2_q09_8.png}{81.2}{0}
\qfigdef{p1_phy_s2_q09_9}{Score-II - Test 34 (02 Feb 2025)-img/p1_phy_s2_q09_9.png}{76.6}{0}
\qfigdef{p1_phy_s2_q10_1}{Score-II - Test 34 (02 Feb 2025)-img/p1_phy_s2_q10_1.png}{59.3}{0}
\qfigdef{p1_phy_s2_q10_2}{Score-II - Test 34 (02 Feb 2025)-img/p1_phy_s2_q10_2.png}{59.3}{0}
\qfigdef{p1_phy_s2_q10_3}{Score-II - Test 34 (02 Feb 2025)-img/p1_phy_s2_q10_3.png}{63.0}{0}
\qfigdef{p1_phy_s2_q10_4}{Score-II - Test 34 (02 Feb 2025)-img/p1_phy_s2_q10_4.png}{91.0}{0}
\qfigdef{p1_phy_s3_q04_1}{Score-II - Test 34 (02 Feb 2025)-img/p1_phy_s3_q04_1.png}{107.3}{0}
\qfigdef{p1_phy_s3_q07_1}{Score-II - Test 34 (02 Feb 2025)-img/p1_phy_s3_q07_1.png}{75.5}{0}
\qfigdef{p1_chem_s1_q01_2}{Score-II - Test 34 (02 Feb 2025)-img/p1_chem_s1_q01_2.png}{24.0}{0}
\qfigdef{p1_chem_s1_q01_1}{Score-II - Test 34 (02 Feb 2025)-img/p1_chem_s1_q01_1.png}{117.2}{0}
\qfigdef{p1_chem_s1_q02_1}{Score-II - Test 34 (02 Feb 2025)-img/p1_chem_s1_q02_1.png}{100.4}{0}
\qfigdef{p1_chem_s1_q02_2}{Score-II - Test 34 (02 Feb 2025)-img/p1_chem_s1_q02_2.png}{18.9}{0}
\qfigdef{p1_chem_s1_q02_4}{Score-II - Test 34 (02 Feb 2025)-img/p1_chem_s1_q02_4.png}{19.0}{0}
\qfigdef{p1_chem_s1_q02_3}{Score-II - Test 34 (02 Feb 2025)-img/p1_chem_s1_q02_3.png}{20.4}{0}
\qfigdef{p1_chem_s1_q02_6}{Score-II - Test 34 (02 Feb 2025)-img/p1_chem_s1_q02_6.png}{26.0}{0}
\qfigdef{p1_chem_s1_q02_5}{Score-II - Test 34 (02 Feb 2025)-img/p1_chem_s1_q02_5.png}{38.8}{0}
\qfigdef{p1_chem_s1_q02_7}{Score-II - Test 34 (02 Feb 2025)-img/p1_chem_s1_q02_7.png}{20.7}{0}
\qfigdef{p1_chem_s1_q04_A}{Score-II - Test 34 (02 Feb 2025)-img/p1_chem_s1_q04_A.png}{79.0}{0}
\qfigdef{p1_chem_s1_q04_B}{Score-II - Test 34 (02 Feb 2025)-img/p1_chem_s1_q04_B.png}{79.1}{0}
\qfigdef{p1_chem_s1_q04_C}{Score-II - Test 34 (02 Feb 2025)-img/p1_chem_s1_q04_C.png}{75.9}{0}
\qfigdef{p1_chem_s1_q04_D}{Score-II - Test 34 (02 Feb 2025)-img/p1_chem_s1_q04_D.png}{75.9}{0}
\qfigdef{p1_chem_s3_q01_1}{Score-II - Test 34 (02 Feb 2025)-img/p1_chem_s3_q01_1.png}{91.6}{0}
\qfigdef{p1_chem_s3_q02_1}{Score-II - Test 34 (02 Feb 2025)-img/p1_chem_s3_q02_1.png}{128.8}{0}
\qfigdef{p1_chem_s3_q03_1}{Score-II - Test 34 (02 Feb 2025)-img/p1_chem_s3_q03_1.png}{192.0}{0}
\qfigdef{p1_chem_s3_q08_1}{Score-II - Test 34 (02 Feb 2025)-img/p1_chem_s3_q08_1.png}{116.5}{0}
\raggedcolumns
\makeatletter\@beginparpenalty=10000\makeatother
\begin{document}
\begin{center}{\Large\bfseries JEE (Main+Advanced) Score-II --- Test 34 (02 Feb 2025)}\end{center}
\begin{multicols}{2}

\partsec{1}{PHYSICS}
\sect{Section-I(i)}

\qn{1} Two identical blocks are placed on a smooth horizontal surface, connected by light string of length $2\ell$. String touches a fixed smooth pulley at its mid-point initially. Shaded parts are two smooth vertical walls. Block A is given a speed $V_0$ perpendicular to string as shown in diagram. B strikes the pulley and stops :- (Pulley and blocks are small)
\qfig{p1_phy_s1_q01_1}
\begin{qoptions}
\opt{A}{Speed of block B when it hits the pulley is $V_0\sqrt{\dfrac38}$}
\opt{B}{Speed of block B when it hits the pulley is $V_0$}
\opt{C}{Speed of A when B hits the pulley is $\sqrt{\dfrac58}V_0$.}
\opt{D}{Speed of A when it hits the wall is $\dfrac{V_0}{2}$}
\end{qoptions}

\qn{2} Density $(\rho)$ versus temperature (T) graph of a thermodynamic cycle of an ideal gas is as shown. If BC and AD are the part of rectangular hyperbola then which of the following graphs will represent the same thermodynamic cycle ?
\qfig{p1_phy_s1_q02_1}
\begin{qoptions}
\opt{A}{\qfigi{p1_phy_s1_q02_A}}
\opt{B}{\qfigi{p1_phy_s1_q02_B}}
\opt{C}{\qfigi{p1_phy_s1_q02_C}}
\opt{D}{\qfigi{p1_phy_s1_q02_D}}
\end{qoptions}

\qn{3} A mass of 0.2 kg is attached to the lower end of a massless spring of force-constant 200 N/m, the upper end of which is fixed to a rigid support. Which of the following statements is/are true ?
\begin{qoptions}
\opt{A}{In equilibrium, the spring will be stretched by 1 cm.}
\opt{B}{If the mass is raised till the spring is unstretched state and then released, it will go down by 2 cm before moving upwards.}
\opt{C}{The frequency of oscillation will be nearly 5 Hz.}
\opt{D}{If the system is taken to the moon, the frequency of oscillation will be the same as on the earth.}
\end{qoptions}

\qn{4} In the potentiometer circuit of given figure the galvanometer reveals a current in the direction shown wherever the sliding contact touches the wire. This could be caused by
\qfig{p1_phy_s1_q04_1}
\begin{qoptions}
\opt{A}{$E_1$ being too low}
\opt{B}{r being too high}
\opt{C}{a break in circuit between PQ}
\opt{D}{$E_2$ being too low}
\end{qoptions}

\qn{5} If electron of the hydrogen atom is replaced by another particle of same charge but of double the mass, then:
\begin{qoptions}
\opt{A}{Bohr radius will increase}
\opt{B}{ionisation energy of the atom will be doubled}
\opt{C}{speed of the new particle in a given state will be lesser than the electron's speed in same orbit}
\opt{D}{gap between energy levels will now be doubled}
\end{qoptions}

\qn{6} Net Electric field due to charges $Q_1$ \& $Q_2$ at a point is E. What will happen to magnitude of electric field if their positions are interchanged.
\begin{qoptions}
\opt{A}{may be increased}
\opt{B}{may be decreased}
\opt{C}{Will remain constant}
\opt{D}{none of these}
\end{qoptions}

\sect{Section-I(ii)}

\qn{7} A spherical interface whose radius is R = 20 cm separates the two media, I and II as shown in the figure. The refractive index of medium I is $\mu_1$ and that of medium II is $\mu_2$. Now match the list$-$I with list$-$II\\
\relax[Consider the pole `O' as the origin and principal axis as the X-axis]
\qfig{p1_phy_s2_q07_1}\par\nobreak
\begin{center}\small\setlength{\tabcolsep}{3pt}
\begin{tabular}{|c|p{3.9cm}|c|p{1.9cm}|}
\hline
 & \multicolumn{1}{c|}{\textbf{List$-$I}} & & \multicolumn{1}{c|}{\textbf{List$-$II}} \\
\hline
(I) & $\mu_2=2,\ \mu_1=1$, object distance $u=-10$ cm & (P) & $v=-40$ cm \\
\hline
(II) & $\mu_2=2,\ \mu_1=1$, object distance $u=-20$ cm & (Q) & $v=-10$ cm \\
\hline
(III) & $\mu_2=1,\ \mu_1=2$, object distance $u=-40$ cm & (R) & $v=40$ cm \\
\hline
(IV) & $\mu_2=1,\ \mu_1=2$, object distance $u=-60$ cm & (S) & $v=-12$ cm \\
\hline
 & & (T) & $v=\infty$ \\
\hline
\end{tabular}
\end{center}
\begin{qoptions}
\opt{A}{I $\to$ T; II $\to$ Q; III $\to$ R; IV $\to$ S}
\opt{B}{I $\to$ P; II $\to$ T; III $\to$ Q; IV $\to$ S}
\opt{C}{I $\to$ P; II $\to$ Q; III $\to$ T; IV $\to$ S}
\opt{D}{I $\to$ Q; II $\to$ P; III $\to$ R; IV $\to$ S}
\end{qoptions}

\qn{8} A uniform rod of length L is suspended from its end by a smooth hinge. The rod is free to rotate in a vertical plane. If the rod is given an angular velocity slightly leaser than $\sqrt{\dfrac{3g}{\ell}}$ at the lowest point, then :-\par\nobreak
\begin{center}\small\setlength{\tabcolsep}{3pt}
\begin{tabular}{|c|p{2.6cm}|c|p{2.6cm}|}
\hline
 & \multicolumn{1}{c|}{\textbf{List-I}} & & \multicolumn{1}{c|}{\textbf{List-II}} \\
\hline
(I) & Force exerted by the hinge in the horizontal direction & (P) & increases continuously till the rod stops \\
\hline
(II) & Force exerted by the hinge in the vertical direction & (Q) & decreases continuously till the rod stops \\
\hline
(III) & Magnitude of Angular velocity of the rod & (R) & increases and then decreases till the rod stops \\
\hline
(IV) & Magnitude of Angular acceleration of the rod & (S) & increases and decreases in a periodic manner \\
\hline
 & & (T) & decreases and increases in a periodic manner \\
\hline
\end{tabular}
\end{center}
\begin{qoptions}
\opt{A}{I $\to$ R,S; II $\to$ Q,T; III $\to$ Q,T; IV $\to$ P,S}
\opt{B}{I $\to$ P,S; II $\to$ Q,T; III $\to$ R,S; IV $\to$ Q,T}
\opt{C}{I $\to$ R,S; II $\to$ Q,S; III $\to$ P,S; IV $\to$ P,T}
\opt{D}{I $\to$ Q,T; II $\to$ R,S; III $\to$ Q,T; IV $\to$ P,S}
\end{qoptions}

\qn{9} The list-I is having the graph of either electrostatic field or its magnitude versus position on x-axis and magnetic field or its magnitude versus position on the x-axis for the system indicated in list-II.\par\nobreak
\begin{center}\small\setlength{\tabcolsep}{2pt}
\begin{tabular}{|c|p{2.0cm}|c|p{2.5cm}|}
\hline
 & \multicolumn{1}{c|}{\textbf{List-I}} & & \multicolumn{1}{c|}{\textbf{List-II}} \\
\hline
(I) & \qfigi{p1_phy_s2_q09_1} & (P) & \qfigi{p1_phy_s2_q09_5}\newline Two infinite wires kept parallel to z-axis carrying current along $+z$-direction. \\
\hline
(II) & \qfigi{p1_phy_s2_q09_2} & (Q) & \qfigi{p1_phy_s2_q09_6}\newline Two infinite wires are kept parallel to z-axis. One carrying current in $+z$-direction, other carrying current in $-z$-direction. \\
\hline
(III) & \qfigi{p1_phy_s2_q09_3} & (R) & \qfigi{p1_phy_s2_q09_7}\newline Two point charges each of magnitude $+Q$ are kept at points $(-a,0)$ and $(a,0)$. \\
\hline
(IV) & \qfigi{p1_phy_s2_q09_4} & (S) & \qfigi{p1_phy_s2_q09_8}\newline Two point charges, $+Q$ and $-Q$ are kept at points $(-a,0)$ and $(a,0)$ respectively \\
\hline
 & & (T) & \qfigi{p1_phy_s2_q09_9}\newline Two infinite non-conducting rods are kept parallel to the y-axis, each carrying uniform charge density $+\lambda$. \\
\hline
\end{tabular}
\end{center}
\begin{qoptions}
\opt{A}{I $\to$ Q,S; II $\to$ P,R,T; III $\to$ Q,S; IV $\to$ P,R,T}
\opt{B}{I $\to$ Q,S; II $\to$ Q,S; III $\to$ P,R,T; IV $\to$ P,R,T}
\opt{C}{I $\to$ Q,S; II $\to$ P,R,T; III $\to$ P,Q,S; IV $\to$ P,R,T}
\opt{D}{I $\to$ P,R,T; II $\to$ Q,S; III $\to$ P,R,T; IV $\to$ P,R,T}
\end{qoptions}

\qn{10} List$-$I contains four different YDSE systems and List$-$II contains intensity measured at some points on the screen. Do the correct match (es) in List$-$I and List$-$II, ($CP_1$ = 0.3 mm and $CP_2$ = 1.2 mm)\par\nobreak
\begin{center}\small\setlength{\tabcolsep}{2pt}
\begin{tabular}{|c|p{4.6cm}|c|p{1.4cm}|}
\hline
 & \multicolumn{1}{c|}{\textbf{List$-$I}} & & \multicolumn{1}{c|}{\textbf{List$-$II}} \\
\hline
(I) & \qfigi{p1_phy_s2_q10_1}\newline Intensity due to each slit $=I_0$ & (P) & $I_C=4I_0$ \\
\hline
(II) & \qfigi{p1_phy_s2_q10_2}\newline Intensity due to each slit $=I_0$ and no absorption of light by water ($\mu_W=4/3$) & (Q) & $I_{P_1}=2I_0$ \\
\hline
(III) & \qfigi{p1_phy_s2_q10_3}\newline Film thickness t = 0.8 $\mu$m, refractive index of film = 3/2\newline Intensity due to each slit $=I_0$ and no absorption of light by glass & (R) & $I_{P_1}=4I_0$ \\
\hline
(IV) & \qfigi{p1_phy_s2_q10_4}\newline Intensity due to each slit $=I_0$ & (S) & $I_{P_1}=0$ \\
\hline
 & & (T) & $I_{P_2}=2I_0$ \\
\hline
\end{tabular}
\end{center}
\begin{qoptions}
\opt{A}{I $\to$ P,Q; II $\to$ P,R; III $\to$ S,T; IV $\to$ P,Q}
\opt{B}{I $\to$ P,R; II $\to$ P,Q; III $\to$ P,Q; IV $\to$ S,T}
\opt{C}{I $\to$ P,R; II $\to$ Q,R; III $\to$ S,T; IV $\to$ P,R}
\opt{D}{I $\to$ P,Q; II $\to$ P,R; III $\to$ P,Q; IV $\to$ S,T}
\end{qoptions}

\sect{Section-II}

\qn{1} A heat engine is operated between two bodies that are kept at constant pressure. The constant pressure heat capacity $C_p$ of the reservoirs is independent of temperature. Initially reservoirs are at temperature 300K and 1200K. If, after some time, they come to a common final temperature $T_f$, the process remaining adiabatic, what is the value of $T_f$ (in Kelvin)?

\qn{2} A small piece of cesium metal (work function `$\phi$' = 1.9 eV) is kept at a distance of 20 cm from a large metal plate having a charge density of $6\epsilon_0$ on surface facing the cesium piece. A monochromatic light of wavelength 400 nm is incident on the cesium piece. Ratio of maximum to minimum kinetic energy of the photoelectron reaching the large metal plate is given by k. Then find the nearest integral value of `k'.

\qn{3} An elastic spherical ball of mass M = 1 kg and radius R = 1 m moving with velocity u = 10 m/s strikes a fixed rigid horizontal surface at an angle $37^\circ$ to the normal. Assuming while it skids in contact with the surface, the coefficient of friction between sphere and the surface is $\mu=0.25$. The angular velocity (in rad/s) of the rebounding ball changes by an amount will be

\qn{4} On four mercury balls lying on the horizontal plane, carefully place the square plate as shown in the figure (top view). The radius of each ball is 1 mm, the mass of the plate is 80 g. There is no wetting. At what distance (in mm) from the horizontal plane will be the lower surface of the plate? Take contact angle to be $180^\circ$.\\
(Take : $\pi=3$, S = 0.49 N/m)
\qfig{p1_phy_s3_q04_1}

\qn{5} The zero of main scale matches with zero of vernier scale. $9^{\text{th}}$ division of main scale matches with $10^{\text{th}}$ division of vernier. No other match is present. After sliding, zero of vernier is found to be in between $3$ \& $4^{\text{th}}$ divisions of main scale \& $5^{\text{th}}$ division of vernier matches with $x^{\text{th}}$ division of main scale. The value of $x/3$ is :

\qn{6} Two conducting plates, each 3 cm $\times$ 6 cm, and three slabs of dielectric, each 1cm $\times$ 3 cm $\times$ 6 cm and having dielectric constants of 1, 2, and 3, are assembled into a parallel plate capacitor with d = 3 cm. Determine the two values of capacitance obtained by the two possible methods of assembling the capacitor. Find $\dfrac{(C_{\text{higher}}-C_{\text{Lower}})}{\epsilon_0}$ (in cm).

\qn{7} A square frame carrying a current I = 0.90 A is located in the same plane as a long straight wire carrying a current $I_0$ = 5A. The frame side has a length a = 8 cm. The axis of the frame passing through the mid-points of the opposite sides is parallel to the wire and is separated from it by a distance h = 15 times greater than the side of the frame. Find the mechanical work (nJ) to be performed in order to turn the frame through $180^\circ$ about its axis, with the current maintained constant. (Take : $\ln\left(\dfrac{31}{29}\right)=0.067$)
\qfig{p1_phy_s3_q07_1}

\qn{8} A small rigid object carries positive and negative charge of magnitude 4 coulomb each. It is oriented so that the positive charge has coordinates $(-1.2\text{ mm},\,1.1\text{ mm})$ and negative charge has coordinates $(1.4\text{ mm},\,-1.3\text{ mm})$. The object is kept in an electric field of $\left(2500\hat i-5000\hat j\right)$ N/C. Find the magnitude of torque (in N-m) acting on the dipole.

\partsec{2}{CHEMISTRY}
\sect{Section-I(i)}

\qn{1} Considering the reaction sequence given below, the correct statement(s) is/are :
\begin{center}\qfigi[0.9]{p1_chem_s1_q01_2}\ \ \qfigi{p1_chem_s1_q01_1}\end{center}
\begin{qoptions}
\opt{A}{Q can give silver mirror image or black ppt with Tollen's reagent.}
\opt{B}{P can give red ppt with Fehling reagent.}
\opt{C}{P and Q can differentiate using Iodoform test.}
\opt{D}{Rate of nucleophilic addition of P greater than Q.}
\end{qoptions}

\qn{2} Consider the following reaction sequence
\qfig{p1_chem_s1_q02_1}
the correct option(s) is/are
\begin{qoptions}
\opt{A}{P = , H$^+$ ; Q = \qfigi{p1_chem_s1_q02_2} ; R = KMnO$_4$/H$^+$}
\opt{B}{P = \qfigi{p1_chem_s1_q02_4}, Alc. KOH, $\Delta$ ; Q = \qfigi{p1_chem_s1_q02_3} ; R = KMnO$_4$/H$^+$}
\opt{C}{P = \qfigi{p1_chem_s1_q02_6}, AlCl$_3$ ; Q = \qfigi{p1_chem_s1_q02_5} ; R = (i) NaOH, I$_2$ (ii) H$_3$O$^+$}
\opt{D}{P = (i) HNO$_3$, H$_2$SO$_4$ (ii) Sn/HCl ; Q = \qfigi{p1_chem_s1_q02_7} ; R = (i) NaNO$_2$ / HCl, (ii) CuCN / KCN (iii) H$_3$O$^+$}
\end{qoptions}

\qn{3} Which of the following is/are \textbf{CORRECT} regarding extraction of iron?
\begin{qoptions}
\opt{A}{Cast iron is made by melting pig iron with scrap iron and coke using hot air blast}
\opt{B}{Wrought iron is purest form of commercial iron \& is prepared from cast iron by oxidizing impurities in reverberatory furnace.}
\opt{C}{Only FeO is produced from Fe$_2$O$_3$ in 500-800 K zone of blast furnace.}
\opt{D}{Pig iron obtained is hard \& brittle and cast iron made from pig iron is soft and ductile.}
\end{qoptions}

\qn{4} Which of the following molecular orbital representation along with probability density function is/are correct for H$_2$?
\begin{qoptions}
\opt{A}{\qfigi{p1_chem_s1_q04_A}}
\opt{B}{\qfigi{p1_chem_s1_q04_B}}
\opt{C}{\qfigi{p1_chem_s1_q04_C}}
\opt{D}{\qfigi{p1_chem_s1_q04_D}}
\end{qoptions}

\qn{5} One litre of 2 molar solution of complex salt CoCl$_3$.4NH$_3$ (mol.wt = 233.5), shows an osmotic pressure of 96 atm at 300 K. The solution is now treated with 2 litre of 1 M AgNO$_3$, which of the following statement(s) is/are correct ? (R = 0.08 L atm/mol K)
\begin{qoptions}
\opt{A}{Mass of AgCl precipitated is 287 g}
\opt{B}{The clear solution will show an osmotic pressure of 96 atm}
\opt{C}{The clear solution will show an osmotic pressure of 32 atm}
\opt{D}{2 moles of [Co(NH$_3$)$_4$Cl$_2$]NO$_3$ will be present in the final solution}
\end{qoptions}

\qn{6} Which of the following statement(s) is/are correct ?
\begin{qoptions}
\opt{A}{It is not possible to compress a gas at a temperature below T$_C$}
\opt{B}{At a temperature below T$_C$, condensation may occur.}
\opt{C}{No condensation takes place above T$_C$}
\opt{D}{Boyle temperature is always greater than T$_C$.}
\end{qoptions}

\sect{Section-I(ii)}

\qn{7} Match the compounds in list-I with the observation in list-II and choose the correct option.\par\nobreak
\begin{center}\small\setlength{\tabcolsep}{3pt}
\begin{tabular}{|c|p{1.9cm}|c|p{3.9cm}|}
\hline
\multicolumn{2}{|c|}{\textbf{List-I}} & \multicolumn{2}{c|}{\textbf{List-II}} \\
\hline
(P) & Glucose & (1) & Deep blue colour (Ruhemann's purple) on addition ninhydring reagent to compound \\
\hline
(Q) & Valine & (2) & Purple/violet coloured complex on addition Molisch's reagent to compound \\
\hline
(R) & Carbolic acid & (3) & Addition of the compound to a saturated solution of NaHCO$_3$ result in effervescence. \\
\hline
(S) & m-toluidine & (4) & The compound reaction with bromine water to give a white precipitate. \\
\hline
 & & (5) & Treating the compound with neutral FeCl$_3$ solution produces violet color. \\
\hline
\end{tabular}
\end{center}
\begin{qoptions}
\opt{A}{P $\to$ 1,2; Q $\to$ 1; R $\to$ 3,5; S $\to$ 4}
\opt{B}{P $\to$ 2; Q $\to$ 1; R $\to$ 3,5; S $\to$ 4}
\opt{C}{P $\to$ 1,2; Q $\to$ 2; R $\to$ 3,5; S $\to$ 4}
\opt{D}{P $\to$ 2; Q $\to$ 1,3; R $\to$ 4,5; S $\to$ 4}
\end{qoptions}

\qn{8}\par\nobreak
\begin{center}\small\setlength{\tabcolsep}{3pt}
\begin{tabular}{|c|p{3.4cm}|c|p{2.1cm}|}
\hline
\multicolumn{2}{|c|}{\textbf{List-I}} & \multicolumn{2}{c|}{\textbf{List-II}} \\
\hline
(P) & $\left[\text{Co(C}_2\text{O}_4)_2(\text{NH}_3)_2\right]^{\ominus}$ & (1) & Low spin complex \\
\hline
(Q) & $\left[\text{RhCl(PPh}_3)_3\right]$ & (2) & Show geometrical isomerism \\
\hline
(R) & $\left[\text{Cr(NH}_3)_2(\text{H}_2\text{O})_4\right]^{+3}$ & (3) & Show optical isomerism \\
\hline
(S) & $\left[\text{Fe(en)}_3\right]^{2+}$ & (4) & Paramagnetic \\
\hline
 & & (5) & Diamagnetic \\
\hline
\end{tabular}
\end{center}
The correct option is :-
\begin{qoptions}
\opt{A}{P $\to$ 1,2,3,5; Q $\to$ 1,4; R $\to$ 1,2,4; S $\to$ 1,4}
\opt{B}{P $\to$ 3,5; Q $\to$ 5; R $\to$ 2,4; S $\to$ 1,3,5}
\opt{C}{P $\to$ 2,5; Q $\to$ 2,3; R $\to$ 1,2,4; S $\to$ 1,2,5}
\opt{D}{P $\to$ 2,3,5; Q $\to$ 2; R $\to$ 4; S $\to$ 1,4}
\end{qoptions}

\qn{9}\par\nobreak
\begin{center}\small\setlength{\tabcolsep}{3pt}
\begin{tabular}{|c|p{2.9cm}|c|p{2.8cm}|}
\hline
\multicolumn{2}{|c|}{\textbf{List-I}} & \multicolumn{2}{c|}{\textbf{List-II}} \\
\hline
(P) & Ammonium\-tetra\-acetato\-plumbate(II) & (1) & FeSO$_4$ + NO $\rightarrow$ \\
\hline
(Q) & Tetra\-ammine\-copper(II) sulphate & (2) & Cu(OH)$_2$.CuSO$_4$ + NH$_3$ $\rightarrow$ \\
\hline
(R) & Iodide of Millon's base & (3) & Cu(OH)$_2$.CuSO$_4$ + K$_4$[Fe(CN)$_6$] $\rightarrow$ \\
\hline
(S) & Nitroso ferrous sulphate & (4) & CH$_3$COONH$_4$ + PbSO$_4$ $\rightarrow$ \\
\hline
 & & (5) & K$_2$[HgI$_4$] + NH$_3$ $\rightarrow$ \\
\hline
\end{tabular}
\end{center}
The correct option is :-
\begin{qoptions}
\opt{A}{P $\to$ 4; Q $\to$ 2; R $\to$ 5; S $\to$ 1}
\opt{B}{P $\to$ 3; Q $\to$ 1; R $\to$ 2; S $\to$ 4}
\opt{C}{P $\to$ 4; Q $\to$ 3; R $\to$ 5; S $\to$ 3}
\opt{D}{P $\to$ 1; Q $\to$ 2; R $\to$ 1; S $\to$ 5}
\end{qoptions}

\qn{10} Following two Lists are provided. The List-I is provided with certain mixtures and List-II is provided with features of resultant colloidal solution.\par\nobreak
\begin{center}\small\setlength{\tabcolsep}{3pt}
\begin{tabular}{|c|p{2.5cm}|c|p{3.3cm}|}
\hline
\multicolumn{2}{|c|}{\textbf{List-I}} & \multicolumn{2}{c|}{\textbf{List-II}} \\
\hline
(P) & AgNO$_3$ + Slight excess KI & (1) & Positively charged sol is formed \\
\hline
(Q) & KI + Slight excess AgNO$_3$ & (2) & Coagulating power is Al$^{3+}>$ Na$^+$ \\
\hline
(R) & FeCl$_3$ + Hot water & (3) & Coagulating value of K$_4$[Fe(CN)$_6$] $<$ NaCl \\
\hline
(S) & FeCl$_3$ + NaOH solution excess & (4) & Negatively charged sol is formed. \\
\hline
 & & (5) & Sol are stable towards addition of electrolytes. \\
\hline
\end{tabular}
\end{center}
Select the option for correct match :
\begin{qoptions}
\opt{A}{P $\to$ 1,3; Q $\to$ 2,4; R $\to$ 2,3; S $\to$ 1,2,5}
\opt{B}{P $\to$ 2,3; Q $\to$ 1,3; R $\to$ 1,2; S $\to$ 1,4}
\opt{C}{P $\to$ 2,4; Q $\to$ 1,4; R $\to$ 2,4; S $\to$ 1,3,5}
\opt{D}{P $\to$ 2,4; Q $\to$ 1,3; R $\to$ 1,3; S $\to$ 2,4}
\end{qoptions}

\sect{Section-II}

\qn{1} Consider the following reaction
\qfig{p1_chem_s3_q01_1}
On estimation of bromine in 1.00 g of R using carius method, the amount of AgBr formed (in g) is \underline{\hspace{1cm}}\\
\relax[Given atomic mass of H = 1, C = 12, O = 16, P = 31, Br = 80, Ag = 108]

\qn{2} The weight percentage of hydrogen in Q formed in the following reaction sequence is
\qfig{p1_chem_s3_q02_1}
\relax[Given : Atomic mass of H = 1, C = 12, N = 14, O = 16, S = 32, Cl = 35]

\qn{3} If the reaction sequence given below is carried out with 4 moles of cyclohexanol, the amount of the product D formed (in g) is \underline{\hspace{1cm}}
\qfig{p1_chem_s3_q03_1}
The yields of A,B,C and D are given in parentheses.\\
\relax[Given : Atomic mass of H = 1, C = 12, O = 16]

\qn{4} Aqueous solution of ammonium chloride on treatment with sodium nitrite yields a major gaseous product ``P'' and trace amount of another gaseous product ``Q''.\\
Find the total number of antibonding electron in P and Q.

\qn{5} Number of specie(s) among the following which have one oxygen/oxide bonded between two similar atoms/ions.\\
Basic beryllium acetate , H$_2$S$_2$O$_8$ , H$_2$S$_2$O$_7$ , N$_2$O$_4$ , N$_2$O$_3$ (symmetrical), $\gamma$-SO$_3$ , H$_4$P$_2$O$_6$ (tetrabasic)

\qn{6} A real gas is subjected to an adiabatic process from (4 bar, 40 L) to (6 bar, 25L) against a constant pressure in a single step. Calculate the enthalpy change for the process (in bar litre).

\qn{7} 1.249 g of a sample containing pure BaCO$_3$ and impure CaCO$_3$ (containing some CaO only) was treated with dil. HCl and it evolved 168 mL of CO$_2$ at 1 atm and 273 K. From this solution, BaCrO$_4$ was precipitated, filtered and washed. The precipitate was dissolved in dilute sulphuric acid and diluted to 100 mL. 10 mL of this solution, when treated with excess KI solution, liberated iodine which required exactly 20 mL of 0.05N Na$_2$S$_2$O$_3$. If `x' is mass \% of CaO in the sample, then `x' will be : (Atomic mass of Ba = 137)\\
(Answer to the nearest integer)

\qn{8} In a 1st order reaction, rate constant varies with temperature according to the graph given below.
\qfig{p1_chem_s3_q08_1}
At 27$^\circ$C, $1.0\times10^{-4}$ \% of the reactant molecules are able to cross the potential barrier. At 72$^\circ$C, the slope of this graph is equal to 0.2 K$^{-1}$ sec$^{-1}$. Calculate the value of rate constant (in s$^{-1}$) at 72$^\circ$C, assuming that activation energy does not change in this temperature range.\\
\relax[Given : R = $\dfrac{25}{3}$ J/K - mol, ln10 = 2.3]

\partsec{3}{MATHEMATICS}
\sect{Section-I(i)}

\qn{1} If $P=2^5.3^6.5^4.7^3$ then number of positive integral divisors of ``P''
\begin{qoptions}
\opt{A}{of form $(2n+3)$, $n\in N$, is = 138}
\opt{B}{of form $(4n+1)$, $n\in W$, is = 70}
\opt{C}{of form $(6n+3)$, $n\in W$, is = 120}
\opt{D}{of form $(4n+3)$, $n\in W$, is = 56}
\end{qoptions}

\qn{2} $f(x)$ and $g(x)$ are two function discontinuous at $x=a$ and $f(x)-g(x)$ is continuous at $x=a$. Also A \& B are symmetric and skew symmetric matrices respectively of order 3, then which of the following statement is/are true
\begin{qoptions}
\opt{A}{$3f(x)+g(x)$ is discontinuous at $x=a$}
\opt{B}{$-f(x)+2g(x)$ may be continuous at $x=a$}
\opt{C}{$AB-BA$ is symmetric}
\opt{D}{$AB^2+B^2A$ is symmetric}
\end{qoptions}

\qn{3} If Bag A has $\Rightarrow$ 3 Red 2 White Balls\\
Bag B has $\Rightarrow$ 2 Red 3 White Balls\\
\& Bag C has $\Rightarrow$ 4 Red 1 White Balls\\
Two balls are drawn at random from any one bag \& found to be both Red then
\begin{qoptions}
\opt{A}{Probability that one more ball drawn from the same bag is white, $=\dfrac12$}
\opt{B}{If one more ball is drawn from the bag C then probability that it is white $=\dfrac7{25}$}
\opt{C}{Probability that the two balls drawn are from bag B $=\dfrac1{10}$}
\opt{D}{Probability that the two balls drawn are from bag A $=\dfrac3{10}$}
\end{qoptions}

\qn{4} If $I_1=\displaystyle\int_0^{\pi/2}\dfrac{\sin(\sin x)}{\sin x}dx,\ I_2=\int_0^{\pi/2}\dfrac{\sin x}{x}dx,\ I_3=\int_0^{\pi/2}\dfrac{\sin(\tan x)}{\tan x}dx$, then which of the following is/are true
\begin{qoptions}
\opt{A}{$I_1>I_3$}
\opt{B}{$I_2>I_3$}
\opt{C}{$I_1>I_2$}
\opt{D}{$I_1>\dfrac\pi2-\dfrac{\pi^3}{144}$}
\end{qoptions}

\qn{5} An ellipse whose major axes is parallel to x-axis is such that the segment of a focal chord are of length 1 and 3 unit. The family of lines $ax+by+c=0$ are the chords of the ellipse such that a,b,c are in A.P. and bisected by the point at which they are concurrent and the equation of auxiliary circle is . $x^2+y^2+2\alpha x+2\beta y-2\alpha-1=0$, then make all the correct alternative(s)
\begin{qoptions}
\opt{A}{The equation of the auxiliary circle $x^2+y^2-2x+4y+1=0$}
\opt{B}{Eccentricity of ellipse is $\dfrac12$}
\opt{C}{Length of major axis is 4}
\opt{D}{Eccentricity of the ellipse is $\dfrac{\sqrt3}2$}
\end{qoptions}

\qn{6} If $x=\dfrac\pi7$ and $y=\dfrac{2\pi}{1999}$, then which of the following is/are true
\begin{qoptions}
\opt{A}{$\cos y.\cos2y.\cos3y\dots\dots\cos999y=\dfrac1{2^{999}}$}
\opt{B}{$\tan^2x+\tan^22x+\tan^23x=21$}
\opt{C}{$8\cos^3x-4\cos^2x-4\cos x=-\dfrac12$}
\opt{D}{$8\cos^3x-4\cos^2x-4\cos x=-1$}
\end{qoptions}

\sect{Section-I(ii)}

\qn{7}\par\nobreak
\begin{center}\small\setlength{\tabcolsep}{3pt}
\begin{tabular}{|c|p{4.1cm}|c|p{1.0cm}|}
\hline
 & \multicolumn{1}{c|}{\textbf{List-I}} & & \textbf{\footnotesize List-II} \\
\hline
(I) & Let $S={}^{31}C_1-\left(1+\tfrac12\right){}^{31}C_2+$\newline $\left(1+\tfrac12+\tfrac13\right){}^{31}C_3-$\newline $\left(1+\tfrac12+\tfrac13+\tfrac14\right){}^{31}C_4+\dots$\newline $+\left(1+\tfrac12+\tfrac13+\dots\right.$\newline $\left.+\tfrac1{31}\right){}^{31}C_{31}$,\newline then $\dfrac1S$ is equal to & (P) & 16 \\
\hline
(II) & If $\dfrac{{}^9C_0}8-\dfrac{{}^9C_1}9+\dfrac{{}^9C_2}{10}-\dfrac{{}^9C_3}{11}+\dots+\dfrac{{}^9C_8}{16}-\dfrac{{}^9C_9}{17}=\dfrac1n$ then n is divisible by & (Q) & 27 \\
\hline
(III) & Remainder when $43^{43^{43}}$ is divided by 40 is & (R) & 13 \\
\hline
(IV) & Let $S_n=\displaystyle\sum_{k=0}^{n}\dfrac{{}^{n+k}C_k}{2^k}$, then remainder when $S_{12}$ is divided by 40 & (S) & 31 \\
\hline
 & & (T) & 17 \\
\hline
\end{tabular}
\end{center}
\begin{qoptions}
\opt{A}{I $\to$ R; II $\to$ R,S; III $\to$ Q; IV $\to$ P}
\opt{B}{I $\to$ S; II $\to$ R,T; III $\to$ Q; IV $\to$ P}
\opt{C}{I $\to$ R; II $\to$ R,Q; III $\to$ S; IV $\to$ S}
\opt{D}{I $\to$ S; II $\to$ R,T; III $\to$ Q; IV $\to$ S}
\end{qoptions}

\qn{8} Let $A=\{1,2,3,4\}$, $B=\{1,2,3,4,5,6,7,8\}$ then match the list I \& II.\par\nobreak
\begin{center}\small\setlength{\tabcolsep}{3pt}
\begin{tabular}{|c|p{4.2cm}|c|p{1.0cm}|}
\hline
 & \multicolumn{1}{c|}{\textbf{List-I}} & & \textbf{\footnotesize List-II} \\
\hline
(I) & Then number of function $f:B\to A$, such that if $x_1>x_2$, then $f(x_1)\ge f(x_2)\ \forall\ x_1\in B$ is N. Then the sum of digits of N is & (P) & 2 \\
\hline
(II) & The number of onto function $f:B\to A$ such that if $x_1>x_2$, then $f(x_1)\le f(x_2)\ \forall\ x_1\in B$ is N. Then the sum of digits of N is & (Q) & 8 \\
\hline
(III) & The number of one-one function $f:A\to B$ such that $f(x)\neq x\ \forall\ x\in A$ is N. Then the sum of digits of N is & (R) & 9 \\
\hline
(IV) & The number of one-one function $f:A\to B$ such that there are exactly two, $x,x\in A$ for which $f(x)=x$ is N. Then the sum of digits of N is & (S) & 7 \\
\hline
 & & (T) & 12 \\
\hline
\end{tabular}
\end{center}
\begin{qoptions}
\opt{A}{I $\to$ Q; II $\to$ T; III $\to$ R; IV $\to$ P}
\opt{B}{I $\to$ Q; II $\to$ T; III $\to$ P; IV $\to$ R}
\opt{C}{I $\to$ T; II $\to$ Q; III $\to$ P; IV $\to$ R}
\opt{D}{I $\to$ T; II $\to$ Q; III $\to$ S; IV $\to$ R}
\end{qoptions}

\qn{9}\par\nobreak
\begin{center}\small\setlength{\tabcolsep}{3pt}
\begin{tabular}{|c|p{4.2cm}|c|p{1.0cm}|}
\hline
 & \multicolumn{1}{c|}{\textbf{List-I}} & & \textbf{\footnotesize List-II} \\
\hline
(I) & Let $f(x)$ be a non negative continuous function $\forall\ x\ge0$, $f(0)=0$. If $f'(x)\le f(x)\sin x\ \ \forall\ x\ge0$, then the value of $f\left(\dfrac{3\pi}2\right)$ is & (P) & 4 \\
\hline
(II) & If $a(x^2+1)+12x+6-3x^2<\sin^{-1}\sin(100)+\cos^{-1}\cos(100)$. $\forall\ x\in R$, $a\in I$, then ``a'' can not be equal to & (Q) & 3 \\
\hline
(III) & If two events A and B are such that $P(A^C)=0.35$, $P(B)=0.4$ and $P(A\cap B)=0.15$, then $20P\left(\dfrac{B}{A\cup B^C}\right)$ is & (R) & 0 \\
\hline
(IV) & The number of elements in the range of the function $f(x)=\sqrt{\cos^{-1}\left(4x^3-3x\right)-\dfrac\pi2}+\dfrac{(2x+3)!}{\sqrt{x+2}}$ is & (S) & 5 \\
\hline
 & & (T) & 8 \\
\hline
\end{tabular}
\end{center}
\begin{qoptions}
\opt{A}{I $\to$ P; II $\to$ P,Q; III $\to$ R; IV $\to$ T}
\opt{B}{I $\to$ R; II $\to$ P,Q,R; III $\to$ P; IV $\to$ T}
\opt{C}{I $\to$ P; II $\to$ Q,R; III $\to$ S; IV $\to$ Q}
\opt{D}{I $\to$ R; II $\to$ P,Q,R,S,; III $\to$ P; IV $\to$ Q}
\end{qoptions}

\qn{10} Four fair dice are rolled simultaneously \& their outcomes are noted. Match the List-I \& List-II\par\nobreak
\begin{center}\small\setlength{\tabcolsep}{3pt}
\begin{tabular}{|c|p{4.1cm}|c|p{1.1cm}|}
\hline
 & \multicolumn{1}{c|}{\textbf{List-I}} & & \textbf{\footnotesize List-II} \\
\hline
(I) & Probability that minimum number obtained on dices is 3 equal to & (P) & $\dfrac{125}{6^4}$ \\
\hline
(II) & Probability that minimum number obtained on dices is 2 and maximum number obtained is 5, is equal to & (Q) & $\dfrac{455}{6^4}$ \\
\hline
(III) & Probability that sum of all scores is 8 & (R) & $\dfrac{175}{6^4}$ \\
\hline
(IV) & Probability that sum of all scores is 16 & (S) & $\dfrac{110}{6^4}$ \\
\hline
 & & (T) & $\dfrac{35}{6^4}$ \\
\hline
\end{tabular}
\end{center}
\begin{qoptions}
\opt{A}{I $\to$ R; II $\to$ S; III $\to$ T; IV $\to$ P}
\opt{B}{I $\to$ R; II $\to$ T; III $\to$ S; IV $\to$ Q}
\opt{C}{I $\to$ R; II $\to$ Q; III $\to$ P; IV $\to$ S}
\opt{D}{I $\to$ P; II $\to$ Q; III $\to$ R; IV $\to$ S}
\end{qoptions}

\sect{Section-II}

\qn{1} Let v and w be distinct randomly chosen roots of the equation $z^{1997}-1=0$. Let $\dfrac mn$ be the probability that $\sqrt{2+\sqrt3}\leqslant|v+w|$, where m and n are relatively prime positive integers. Find $m+n$.

\qn{2} If $P=\dfrac{xyz}{(1+5x)(4x+3y)(5y+6z)(z+18)}$ where $x,y,z\in R^+$ then minimum value of $\dfrac1p$ is equal to

\qn{3} Let A is non singular matrix of order two such that $A+\text{adj}A=A^{-1}$ and $|2A^{-1}|=x$. Let P \& Q are square matrices of same order such that $|P|=|Q|=1$ and $P(\text{adj}\,P+\text{adj}\,Q)=Q$ then value of $|P+Q|=y$ then value of $x+y$ is equal to

\qn{4} Six points $(x_i,y_i)$, $i=1,2,\dots\dots,6$ are taken on the circle $x^2+y^2=4$ such that $\displaystyle\sum_{i=1}^{6}x_i=8$ and $\displaystyle\sum_{i=1}^{6}y_i=4$. The line segment joining orthocentre of a triangle made by any three points and the centroid of the triangle made by other three points passes through a fixed points $(h,k)$, then $h+k$ is

\qn{5} If $f:(0,\infty)\to N$ and $f(x)=\left[\dfrac{x^2+x+1}{x^2+1}\right]+\left[\dfrac{4x^2+x+2}{2x^2+1}\right]+\left[\dfrac{9x^2+x+3}{3x^2+1}\right]+\dots+\left[\dfrac{n^2x^2+x+n}{nx^2+1}\right]$, $n\in N$, then the value of $\displaystyle\lim_{n\to\infty}\left[\dfrac{f(x)-n}{(f(x))^2-\frac{n^3(n+2)}4}\right]=\ell_1$, (where $[y]$ denotes the greatest integer less than or equal to y)\\
Also Let $P(x)=-x^3+x^2-x+1$\\
$g(x)=\begin{cases}\min(P(t)),\ 0\leqslant t\leqslant x,\ 0\leqslant x\leqslant1\\ x-1,\ 1<x\leqslant2\end{cases}$, then $\displaystyle\lim_{x\to1^+}g(g(x))=\ell_2$ then value of $\ell_1+\ell_2$ is

\qn{6} Let $|\vec a|=1$, $|\vec b|=1$, $|\vec a+\vec b|=\sqrt3$ and $\vec c$ be a vector such that $\vec c=\vec a+2\vec b-3(\vec a\times\vec b)$. If $y=|(\vec a\times\vec b)\times\vec c|$ then value of $[y^2]=$ \underline{\hspace{1cm}} (where [.] denotes G.I.F.)

\qn{7} Plane $x+y+z=10$ intersects the x-axis, y-axis and z-axis at points A, B and C respectively then number of points with positive integral coordinates inside tetrahedron OABC is equal to \underline{\hspace{1cm}}

\qn{8} Let P be a variable point in I-octant, such that distance of P from x-axis is $d_1$, from y-axis is $d_2$, and from z-axis is $d_3$. If locus of point P is such that $\sqrt7d_1+\sqrt8d_2+3d_3=12$, then square of minimum distance of P from origin is

\end{multicols}
\end{document}
